% Figure: the local driving temperature difference theta = T_h - T_c along the
% exchanger, decaying as a single exponential in UA for counter-flow and for
% parallel-flow. The figure makes visible the result of the differential balance:
% theta varies geometrically along the area, which is what produces the log-mean.
\begin{figure}[htbp]
\centering
\begin{tikzpicture}
\begin{axis}[
  width=12cm, height=6.5cm,
  xlabel={fraction of area traversed, $A^\ast = A_{\text{partial}}/A$},
  ylabel={driving difference $\theta/\theta_{\text{inlet}}$},
  xmin=0, xmax=1, ymin=0, ymax=1.05,
  axis lines=left,
  legend pos=north east,
  legend cell align=left,
  tick label style={font=\scriptsize},
  label style={font=\small},
  legend style={font=\scriptsize},
  samples=80, domain=0:1,
]
% Counter-flow: theta decays as exp(-NTU(1-Cr) A*) but stays well above zero at exit;
% drawn for NTU=1.5, Cr=0.6 normalised to its larger end value.
\addplot[engblue, very thick] {exp(-1.5*(1-0.6)*x)};
\addlegendentry{counter-flow ($\theta_2/\theta_1$ moderate)}
% Parallel-flow: theta decays as exp(-NTU(1+Cr) A*), collapsing toward a small exit
% value; drawn for NTU=1.5, Cr=0.6.
\addplot[engred, very thick] {exp(-1.5*(1+0.6)*x)};
\addlegendentry{parallel-flow ($\theta_2/\theta_1$ small)}
% the geometric-mean intuition: log-linear decay means the area-average is the log-mean
\node[engblue, font=\scriptsize, anchor=west] at (axis cs:0.55,0.62) {slope $-\mathrm{NTU}(1-C_r)$};
\node[engred, font=\scriptsize, anchor=west] at (axis cs:0.40,0.18) {slope $-\mathrm{NTU}(1+C_r)$};
\end{axis}
\end{tikzpicture}
\caption[Driving difference decaying along the exchanger]{Local driving
temperature difference $\theta = T_h - T_c$ along the exchanger, normalised to its
value at the inlet end, for a counter-flow and a parallel-flow arrangement at the
same number of transfer units and the same capacity ratio. The differential
balance $d\theta/\theta = -(U/C^\ast)\,dA$ integrates to a single exponential, so
$\theta$ falls log-linearly along the area and the appropriate area-average of a
log-linearly varying quantity is its log-mean, which is the origin of the log-mean
temperature difference. The counter-flow curve decays with the gentler slope
$-\mathrm{NTU}(1-C_r)$ and keeps a usable driving difference all the way to the
exit, while the parallel-flow curve decays with the steeper slope
$-\mathrm{NTU}(1+C_r)$ and collapses toward the common outlet temperature the two
streams must share. The larger area under the counter-flow curve is the same fact,
read geometrically, as its larger log-mean difference for fixed terminal
temperatures.}
\label{fig:vol04ch06:deltat-decay}
\end{figure}
